\documentclass[a4paper,12pt]{article}

\usepackage{iftex}
\ifPDFTeX
  \usepackage[T2A]{fontenc}
  \usepackage[utf8]{inputenc}
  \usepackage[russian]{babel}
  \usepackage{helvet}
  \renewcommand{\familydefault}{\sfdefault}
\else
  \usepackage{fontspec}
  \usepackage{polyglossia}
  \setdefaultlanguage{russian}
  \setmainfont{DejaVu Sans}
  \setsansfont{DejaVu Sans}
  \renewcommand{\familydefault}{\sfdefault}
\fi

\usepackage{amsmath,amssymb,mathtools}
\usepackage{icomma}
\usepackage{geometry}
\usepackage{microtype}
\usepackage{tikz}
\usepackage{tabularx,booktabs,longtable,array}
\usepackage{enumitem}
\usepackage{hyperref}
\usepackage{parskip}
\usepackage{fancyhdr}
\usepackage{setspace}
\usepackage{xcolor}

\geometry{left=18mm,right=18mm,top=18mm,bottom=20mm}
\onehalfspacing
\setlength{\parindent}{0pt}
\setlength{\parskip}{6pt}
\setlength{\emergencystretch}{3em}
\setlist[itemize]{leftmargin=7mm,itemsep=3pt,topsep=4pt}
\setlist[enumerate]{leftmargin=8mm,itemsep=4pt,topsep=4pt}
\renewcommand{\arraystretch}{1.35}

\definecolor{accent}{HTML}{7A3B45}
\definecolor{darkaccent}{HTML}{4D2730}
\definecolor{textgray}{HTML}{333333}
\definecolor{softgray}{HTML}{6E6E6E}
\definecolor{paperline}{HTML}{D8C8BC}
\definecolor{success}{HTML}{286B5A}

\hypersetup{
  colorlinks=true,
  linkcolor=darkaccent,
  urlcolor=accent,
  pdftitle={Финансовая математика: сравнение вкладов и поиск ставки},
  pdfauthor={Лёвин Артём Александрович}
}
\pagestyle{fancy}
\fancyhf{}
\lhead{\textcolor{textgray}{Финансовая математика}}
\rhead{\textcolor{textgray}{Вклады и сравнение сценариев}}
\cfoot{\textcolor{softgray}{\thepage}}
\setlength{\headheight}{15pt}

\newcommand{\term}[1]{\textcolor{darkaccent}{\textbf{#1}}}
\newcommand{\formula}[1]{%
  \begin{center}
    \begin{minipage}{0.94\linewidth}
      \centering\large\color{darkaccent}#1
    \end{minipage}
  \end{center}}
\newcommand{\checkitem}{\item[$\square$]}
\newcommand{\rub}[1]{\,\text{руб.}}
\newcommand{\thrub}[1]{\,\text{тыс. руб.}}

\begin{document}

\begin{titlepage}
  \thispagestyle{empty}
  \centering
  \vspace*{25mm}
  {\Large\textcolor{accent}{Чек-лист}\par}
  \vspace{7mm}
  {\Huge\bfseries\textcolor{darkaccent}{Финансовая математика}\par}
  \vspace{3mm}
  {\LARGE Сравнение вкладов и поиск процентной ставки\par}
  \vfill
  {\large Лёвин Артём Александрович\par}
  {\large эксклюзивно для Николь Саркисянц\par}
  \vspace{8mm}
  {\large \today\par}
  \vfill
  \begin{tikzpicture}[x=1cm,y=1cm]
    \draw[accent,line width=1.2pt]
      (-7,0) .. controls (-4.8,1.2) and (-2.4,-1.1) .. (0,0)
             .. controls (2.4,1.1) and (4.8,-1.2) .. (7,0);
    \draw[paperline,line width=.6pt]
      (-7,-.35) .. controls (-3.5,.45) and (3.5,-.45) .. (7,-.35);
  \end{tikzpicture}
  \vspace*{8mm}
\end{titlepage}

\tableofcontents
\clearpage

\section{Главная идея урока}

Задачи о вкладах удобно решать не длинной цепочкой денежных вычислений, а через
\term{коэффициент роста} и рекуррентную модель. Такой способ сохраняет структуру
задачи, помогает сравнивать сценарии и часто полностью исключает неизвестную
первоначальную сумму.

\begin{table}[ht]
\centering
\caption{Единый язык финансовой модели}
\begin{tabularx}{\linewidth}{@{}>{\bfseries}lX@{}}
\toprule
Обозначение & Смысл \\
\midrule
$S$ & первоначальная сумма вклада; \\
$P\%$ & процентная ставка за один год; \\
$p=\dfrac{P}{100}$ & ставка в виде десятичной дроби; \\
$r=1+p=1+\dfrac{P}{100}$ & коэффициент роста за один год; \\
$A_k$ & операция после начисления процентов в конце $k$-го года:
$A_k>0$ --- взнос, $A_k<0$ --- снятие; \\
$S_k$ & сумма после начисления и операции в конце $k$-го года; \\
$\Delta S$ & разность сравниваемых сумм, обычно
$S_{\text{план}}-S_{\text{факт}}$. \\
\bottomrule
\end{tabularx}
\end{table}

\formula{$\boxed{S_k=rS_{k-1}+A_k}$}

\term{Порядок событий:}
\[
\text{остаток в начале года}
\longrightarrow
\text{начисление процентов}
\longrightarrow
\text{взнос или снятие}.
\]
Если операция совершается после начисления, переносить её перед умножением на
$r$ нельзя.

\section{Универсальный алгоритм}

\begin{enumerate}
  \item Обозначьте первоначальную сумму буквой $S$.
  \item Переведите ставку: $p=P/100$ и запишите $r=1+p$.
  \item Зафиксируйте знаки операций $A_k$: взнос со знаком «плюс», снятие со знаком «минус».
  \item Проверьте порядок событий в каждом году.
  \item Запишите плановый сценарий без операций: после трёх лет это $Sr^3$.
  \item Постройте фактический сценарий рекуррентно или таблицей.
  \item Составьте нужную разность либо неравенство и раскрывайте скобки поэтапно.
  \item Решите полученное уравнение, проверьте все корни по смыслу, единицы и округление.
\end{enumerate}

\begin{table}[ht]
\centering
\caption{Трёхлетняя модель с операциями после начисления}
\begin{tabularx}{\linewidth}{@{}cXXX@{}}
\toprule
Год & Начало года & Операция & Конец года \\
\midrule
1 & $S$ & $A_1$ & $Sr+A_1$ \\
2 & $Sr+A_1$ & $A_2$ & $(Sr+A_1)r+A_2$ \\
3 & $(Sr+A_1)r+A_2$ & $A_3$ & $\bigl((Sr+A_1)r+A_2\bigr)r+A_3$ \\
\bottomrule
\end{tabularx}
\end{table}

Если $A_3=0$, то
\[
S_3=Sr^3+A_1r^2+A_2r.
\]
Слагаемое $A_1$ работает два последующих года, поэтому умножается на $r^2$;
слагаемое $A_2$ работает один год, поэтому умножается на $r$.

\section{Разобранные задачи}

\subsection{Задача 1. Снятие и последующее пополнение}

Вклад открыт под $8\%$ годовых. После первого начисления сняли
$3000$ рублей, после второго внесли $3000$ рублей. Найдём, на сколько
фактическая сумма через три года меньше плановой.

\[
p=0{,}08,\qquad r=1{,}08,\qquad
A_1=-3000,\quad A_2=3000,\quad A_3=0.
\]

\term{План:}
\[
S_{\text{план}}=Sr^3.
\]

\term{Факт:}
\begin{align*}
S_1&=Sr-3000,\\
S_2&=(Sr-3000)r+3000,\\
S_{\text{факт}}=S_3&=\bigl((Sr-3000)r+3000\bigr)r\\
&=Sr^3-3000r^2+3000r.
\end{align*}

\term{Сравнение:}
\begin{align*}
\Delta S
&=S_{\text{план}}-S_{\text{факт}}\\
&=Sr^3-\bigl(Sr^3-3000r^2+3000r\bigr)\\
&=3000r^2-3000r\\
&=3000r(r-1)\\
&=3000\cdot1{,}08\cdot0{,}08=259{,}20.
\end{align*}

\term{Ответ:} фактическая сумма меньше плановой на $259{,}20$ рубля.

\term{Смысл результата.} Снятые $3000$ рублей не приносили проценты два года,
а возвращённые $3000$ рублей работали только последний год. Поэтому одинаковые
по модулю операции не компенсируют друг друга полностью.

\subsection{Задача 2. Неизвестная целая ставка}

После первого начисления со вклада сняли $20000$ рублей, после второго внесли
$10000$ рублей, после третьего операций не было. Через три года фактическая
сумма оказалась на $16800$ рублей меньше плановой. Ставка --- целое число
процентов. Найдём её.

\[
A_1=-20000,\qquad A_2=10000,\qquad A_3=0.
\]
Фактическая сумма равна
\[
S_{\text{факт}}=Sr^3-20000r^2+10000r.
\]
Следовательно,
\begin{align*}
16800
&=S_{\text{план}}-S_{\text{факт}}\\
&=Sr^3-\bigl(Sr^3-20000r^2+10000r\bigr)\\
&=20000r^2-10000r.
\end{align*}

Получаем квадратное уравнение:
\[
20000r^2-10000r-16800=0,
\]
\[
50r^2-25r-42=0.
\]
\[
D=(-25)^2-4\cdot50\cdot(-42)=625+8400=9025=95^2.
\]
\[
r_{1,2}=\frac{25\pm95}{100},\qquad
r_1=1{,}2,\qquad r_2=-0{,}7.
\]

Оба корня квадратного уравнения вычислены, но $r=-0{,}7$ не подходит:
коэффициент роста вклада не может быть отрицательным. Поэтому
\[
r=1{,}2,\qquad p=r-1=0{,}2,\qquad P=100p=20.
\]

\term{Ответ:} $20\%$.

\subsection{Задача 3. Фиксированное ежегодное пополнение}

Владимир положил в банк $3600$ тысяч рублей под $10\%$ годовых. В конце
каждого из первых двух лет после начисления процентов он вносил одну и ту же
сумму $A$. К концу третьего года вклад оказался на $48{,}5\%$ больше
первоначального. Найдём $A$.

Фраза «больше на $48{,}5\%$» означает:
\[
S_3=S\left(1+\frac{48{,}5}{100}\right)=1{,}485S.
\]
При $r=1{,}1$:
\[
S_3=Sr^3+Ar^2+Ar.
\]
Значит,
\[
1{,}331S+A(1{,}21+1{,}1)=1{,}485S,
\]
\[
2{,}31A=0{,}154S.
\]
Подставляем $S=3600$ тыс. руб.:
\[
A=\frac{0{,}154\cdot3600}{2{,}31}=240\text{ тыс. руб.}
\]

\term{Ответ:} $240$ тысяч рублей.

\subsection{Задача 4. Сравнение двух вкладов}

По вкладу $A$ банк три года подряд начисляет по $10\%$. По вкладу $B$
первые два года начисляется по $11\%$, а в третий год --- целое число $x\%$.
Найдём наименьшее $x$, при котором вклад $B$ за три года окажется выгоднее.

\[
S_A=S\cdot1{,}1^3=1{,}331S,
\]
\[
S_B=S\cdot1{,}11^2\left(1+\frac{x}{100}\right)
=1{,}2321S\left(1+\frac{x}{100}\right).
\]
Требуется строгое неравенство $S_B>S_A$. Так как $S>0$, сокращаем на $S$:
\[
1{,}2321\left(1+\frac{x}{100}\right)>1{,}331.
\]
\[
\frac{x}{100}>
\frac{1{,}331-1{,}2321}{1{,}2321}
=\frac{989}{12321}\approx0{,}080269.
\]
\[
x>8{,}0269.
\]
Наименьшее целое значение, удовлетворяющее строгому неравенству, равно $9$.

\term{Ответ:} $9\%$.

\section{Особые случаи и смысловые проверки}

\subsection{Как выбрать порядок вычитания}

Если сказано «фактическая сумма меньше плановой на $d$», то
\[
S_{\text{план}}-S_{\text{факт}}=d.
\]
Если направление сравнения неизвестно, допустимо сначала записать
\[
\left|S_{\text{план}}-S_{\text{факт}}\right|=d,
\]
а затем проверить полученные варианты по смыслу.

\subsection{Почему сокращается первоначальная сумма}

В разности планового и фактического сценариев одинаковые слагаемые $Sr^3$
уничтожаются. Это закономерно: разница вызвана только дополнительными
операциями. Если после раскрытия скобок неизвестная сумма $S$ не сократилась
в задаче, где она не дана и не требуется, следует проверить модель.

\subsection{Почему нельзя сразу отбрасывать второй корень}

Квадратное уравнение нужно решить полностью. Лишь после вычисления корней
применяется смысловая область:
\[
r=1+\frac{P}{100}>0.
\]
При положительной процентной ставке дополнительно $r>1$. Отрицательный корень
не «исчезает», а отвергается с обоснованием.

\subsection{Строгое сравнение и округление}

При словах «выгоднее», «больше» используется знак $>$. Если найдено
$x>8{,}0269$ и $x$ должен быть целым, ответ --- $9$, а не $8$.
Округление выполняется не механически, а так, чтобы условие осталось истинным.

\section{Типичные ошибки}

\begin{enumerate}
  \item Использовать $P$ вместо $p=P/100$ и записывать $r=1+P$.
  \item Поменять порядок действий: сначала выполнить операцию, затем начислить проценты.
  \item Потерять знак минус перед скобкой при вычислении
  $S_{\text{план}}-S_{\text{факт}}$.
  \item Раскрыть все скобки одной строкой и смешать степени $r$ с денежными суммами.
  \item Решить квадратное уравнение, но не проверить корни по смыслу.
  \item Перевести «больше на $q\%$» как $A=B+q$ вместо
  $A=B(1+q/100)$.
  \item Сравнить суммы, относящиеся к разным моментам времени.
  \item Округлить вниз при строгом условии «выгоднее».
\end{enumerate}

\clearpage
\section{Итоговый чек-лист}

Перед записью ответа проверьте:

\begin{itemize}
  \checkitem Я обозначил(а) первоначальную сумму и процентную ставку.
  \checkitem Я перевёл(а) проценты в десятичную дробь.
  \checkitem Я записал(а) коэффициент роста $r=1+P/100$.
  \checkitem Я отметил(а) каждый взнос знаком «плюс», а снятие знаком «минус».
  \checkitem Я учёл(а), что операция выполняется после начисления процентов.
  \checkitem Я записал(а) рекуррентную формулу $S_k=rS_{k-1}+A_k$.
  \checkitem Я сравниваю плановую и фактическую суммы в один момент времени.
  \checkitem Я раскрываю вложенные скобки поэтапно и отдельно проверяю знаки.
  \checkitem Я решил(а) уравнение полностью и обосновал(а) отбор корней.
  \checkitem Я перевёл(а) найденный коэффициент $r$ обратно в проценты.
  \checkitem Я учёл(а) строгость неравенства и округлил(а) по смыслу.
  \checkitem Я проверил(а) единицы измерения и здравый смысл ответа.
\end{itemize}

\clearpage
\section*{Тренировочный блок}
\addcontentsline{toc}{section}{Тренировочный блок}
\begin{spacing}{1.22}

\textbf{Средний уровень}

\begin{enumerate}
  \item Вклад $60000$ рублей открыт под $5\%$ годовых. Найдите сумму через
  два года, если дополнительных операций не было.

  \item На вклад положили $100000$ рублей под $10\%$ годовых. После первого
  начисления внесли ещё $5000$ рублей. После второго начисления операций не
  было. Найдите итоговую сумму.

  \item На вклад положили $90000$ рублей под $12\%$ годовых. После первого
  начисления сняли $15000$ рублей. Найдите сумму после второго начисления.
\end{enumerate}

\textbf{Уровень выше среднего}

\begin{enumerate}[resume]
  \item Вклад открыт под $8\%$ годовых. После первого начисления сняли
  $3000$ рублей, после второго внесли $3000$ рублей. На сколько фактическая
  сумма через три года меньше плановой?

  \item Владимир положил в банк $3600$ тысяч рублей под $10\%$ годовых. После
  начисления процентов в конце каждого из первых двух лет он вносил одну и ту
  же сумму. К концу третьего года вклад стал на $48{,}5\%$ больше
  первоначального. Найдите ежегодный взнос.

  \item Вклад $A$ три года растёт на $10\%$ ежегодно. Вклад $B$ первые два
  года растёт на $11\%$, а за третий год --- на целое число процентов.
  Найдите наименьшую ставку третьего года, при которой вклад $B$ выгоднее.

  \item Ставка --- целое число процентов. После первого начисления сняли
  $12000$ рублей, после второго внесли $6000$ рублей. Через три года
  фактическая сумма оказалась на $8970$ рублей меньше плановой. Найдите ставку.

  \item На вклад положили $250000$ рублей под $7\%$ годовых. После первого
  начисления внесли $10000$ рублей, после второго сняли $5000$ рублей,
  после третьего операций не было. Найдите итоговую сумму.

  \item Вклад $A$ три года растёт на $9\%$ ежегодно. Вклад $B$ первые два
  года растёт на $10\%$, а за третий год --- на целое число процентов.
  Найдите наименьшую ставку третьего года, при которой вклад $B$ выгоднее.
\end{enumerate}

\textbf{Сложный уровень}

\begin{enumerate}[resume]
  \item Вклад открыт под $10\%$ годовых. После первого начисления сняли
  $x$ рублей, после второго внесли $0{,}4x$ рублей, после третьего операций
  не было. Через три года фактическая сумма оказалась на $7700$ рублей меньше
  плановой. Найдите $x$.
\end{enumerate}
\end{spacing}

\clearpage
\section*{Ответы}
\addcontentsline{toc}{section}{Ответы}

\begin{center}
\textcolor{softgray}{Сначала выполните все задания самостоятельно.}\\[4pt]
\textbf{Код самопроверки: 1807}
\end{center}

\begin{longtable}{@{}cl@{}}
\toprule
\textbf{№} & \textbf{Ответ} \\
\midrule
\endhead
1 & $66150$ руб. \\
2 & $126500$ руб. \\
3 & $96096$ руб. \\
4 & $259{,}20$ руб. \\
5 & $240$ тыс. руб. \\
6 & $9\%$ \\
7 & $15\%$ \\
8 & $312359{,}75$ руб. \\
9 & $8\%$ \\
10 & $10000$ руб. \\
\bottomrule
\end{longtable}

\end{document}
